\documentclass[a4paper, 12pt]{article}
\usepackage[top=3cm, bottom=3cm, left = 2.5cm, right = 2cm]{geometry}
\geometry{a4paper}
\usepackage{booktabs}
\usepackage{multirow}
\usepackage{array}
\usepackage{threeparttable}
\usepackage{makecell}
\usepackage{adjustbox}
\usepackage{amsmath}
\usepackage{hyperref}

\usepackage{url}
\usepackage[utf8]{inputenc}
% T1 + lmodern give the typewriter font a real em dash and real straight
% quotes; without them the verbatim prompts in the appendix render "—" as a
% vertical bar and "'" as a curly quote, i.e. not verbatim at all.
\usepackage[T1]{fontenc}
\usepackage{lmodern}
\usepackage{upquote}
\usepackage{newunicodechar}
\newunicodechar{≤}{\ensuremath{\le}}
\newunicodechar{≥}{\ensuremath{\ge}}

% Verbatim prompt blocks for the appendix. `literate` maps the only non-ASCII
% character the prompts contain (U+2014 em dash) so the listing stays faithful.
\usepackage{listings}
\lstdefinestyle{prompt}{
  basicstyle=\ttfamily\scriptsize,
  breaklines=true,
  breakautoindent=false,
  breakindent=0pt,
  columns=fullflexible,
  keepspaces=true,
  frame=single,
  framesep=4pt,
  xleftmargin=2pt, xrightmargin=2pt,
  aboveskip=5pt, belowskip=5pt,
  literate={—}{{\textemdash}}1
}
\lstset{style=prompt}

% Tighter float and section spacing
\usepackage[skip=4pt]{caption}
\usepackage{titlesec}
\titlespacing*{\section}{0pt}{1.6ex plus .2ex}{0.8ex}
\titlespacing*{\subsection}{0pt}{1.3ex plus .2ex}{0.5ex}
\setlength{\textfloatsep}{8pt}
\setlength{\intextsep}{8pt}

\title{What Reflection Is Told, Not Where It Sits:\\
A Controlled Study of Self-Reflection in\\
Multi-Turn Tool Calling}
\author{Mirodil Kamilov}
\date{SS 2026}

\begin{document}
\maketitle

\begin{abstract}
\noindent
Self-reflection is a widely adopted way to improve LLM agents: an agent whose
episode fails writes down what went wrong, and that lesson is carried into a
fresh attempt. Its largest reported gains come from settings with rich feedback,
such as program synthesis. This report asks whether it survives a
harder setting --- BFCL~v3 multi-turn, where an agent sustains a stateful
conversation across several turns of API calls and the feedback a reflection
step may legitimately receive is far thinner. I evaluate Reflexion against a
\emph{sham} reflection, matched in form and length to a real one but carrying
only general advice, never a lesson from the failure. The sham performed slightly better
($-1.5$~pp for reflection, 95\% CI $[-5.5,+2.5]$). Supplying the reflection step
with a substantially richer input --- a record of how the agent's own actions changed
the environment on each turn --- also changed nothing ($+0.0$~pp). Reflection is nevertheless
not inert. Across six conditions of increasing signal quality, only 28 of 115 failed tasks
are recovered by \emph{any} retry, and every condition recovers 15--18 of those same 28:
reflection alters \emph{which} failures are recovered rather than \emph{how many}.
\end{abstract}

\section{Introduction}

An LLM agent solves a task by alternating between reasoning and acting: it
inspects its situation, calls a tool, reads the result, and repeats. Two
influential proposals add an explicit reflective step to that loop. ReAct
\cite{yao2023react} has the model write a short thought before each action.
Reflexion \cite{shinn2023reflexion} waits until an episode has failed, has the
model write a verbal lesson explaining the failure, and carries that lesson into
a fresh attempt. Between them sits a third position --- reflecting after each
individual tool result, to decide whether the last action should be revised,
which the REBACT \cite{zeng2025rebact} line of work explores.

I set out to compare these three positions directly. The plan was to build five
architectures --- a plain function-calling baseline, ReAct, episodic Reflexion,
a Reflexion variant with a persistent cross-task memory, and REBACT --- run them
over one frozen set of tasks, and answer a single question: \emph{where in the
reason--act loop does reflection actually help?}

\subsection{Why this benchmark, and why the plan did not survive it}

The original ReAct and Reflexion results were obtained largely on ALFWorld,
HotpotQA and WebShop. Those tasks are demanding, and some are long — an ALFWorld
episode can run for tens of steps — but they share a narrow notion of what the agent
must get right: each episode carries a single instruction, chosen from a small action vocabulary,
and is scored on whether the goal was reached. HotpotQA, for instance, gives the agent three actions.

BFCL~v3 \cite{patil2025bfcl} is however a \emph{multi-turn} and \emph{multi-step} benchmark: an episode
consists of several realistic user–agent interactions, and each turn may require chaining several calls
to satisfy one request. A task in my subset exposes a median of 28 callable functions (range 17--39).
The agent is graded on both the resulting environment state and the minimum set of calls the task requires.
A further part of the difficulty is about not acting at all: two categories test abstention,
which neither earlier benchmark does at all — whether the agent asks for an argument
the user omitted (\texttt{miss\_param}), and whether it recognises that no available tool fits (\texttt{miss\_func}).

A final difference bears on reflection directly. Reflexion's largest gains come from program synthesis,
where the reflection step reads a failing test suite — a signal close to a diagnosis in itself.
BFCL offers no test output and no partial credit, only a verdict, so a reflection step must infer its own diagnosis.

Three findings, in order, made the original comparison untenable.

\textbf{First, the measurement floor turned out to be high.} The model used here,
\texttt{qwen3-next-80b-a3b-instruct} \cite{qwen3}, is a mixture-of-experts model
served over a shared endpoint, and at \texttt{temperature=0} its outputs are still
not reproducible. Running the identical baseline twice over the same 200 tasks produced
different verdicts on 15 of them. That places a floor on what the experiment can resolve: any true
effect below roughly 5 percentage points (pp) is indistinguishable from noise,
and the effects under investigation are plausibly that size.

\textbf{Second, the ReAct condition produced a null for a reason specific to
modern models.} ReAct was proposed for base models that do not spontaneously
verbalise a plan, so instructing them to do so supplies something genuinely
absent. The instruction-tuned model used here already writes free-text reasoning
before 80\% of its tool calls without any prompting. The scaffold therefore had
very little left to contribute, and the pass rate did not move. The result
characterises the model rather than the method.

\textbf{Third, and decisively, a control arm behaved unexpectedly.} To separate
the value of a reflection's \emph{content} from the value of simply attempting
the task again, I added an arm in which the model's own lesson was replaced by a
fixed paragraph, matched in form and length to a real reflection but saying
nothing about the actual failure --- a \emph{sham}. The expectation was that
this arm would fall clearly behind genuine reflection. It did not; it scored
marginally higher. Reducing the sham's lesson further, to a single bland sentence,
changed nothing again. Whatever was producing the gain, it was not the content
of the reflection.

That observation is not isolated: Kapoor et al. \cite{kapoor2024agents} find that on
HumanEval a bare retry baseline — re-invoking the model at temperature zero whenever
the tests fail — matches the best published agent architectures, precisely because such models
are not deterministic even at temperature zero. Their comparison runs between architectures; this report
isolates the same effect within one, by holding the retry fixed and varying only what the reflection is told.

\subsection{The revised question}

If reflection was contributing nothing, one natural explanation was that it had
too little to work with: the agent was told only \emph{which} simulated API
ended in the wrong state, or on \emph{which} turn it failed. I therefore built a
further arm that additionally shows the reflection step a record of what its own
actions changed in the environment on each turn --- the difference between the
environment state before and after every turn of its failed attempt. This is
more information, it is cheap, and it cannot leak the expected answer. It too
changed nothing.

The question consequently shifted from \emph{where} reflection sits in the loop
to \emph{what it is told}:

\begin{quote}
Does the information content of the reflection signal change \emph{how many}
failures an agent recovers, or only \emph{which ones}?
\end{quote}

The answer is: only which ones. This seminar project's contribution is adding the
evaluation --- six-condition ladder of steadily richer reflection input over one frozen task set
and a sham control that separates reflection from retrying (Section~\ref{subsec:ladder}); evidence that
recoverable failures form a small fixed pool, so reflection reallocates rather
than enlarges the set of tasks solved (Sections~\ref{sec:which}--\ref{sec:ceiling});
and a hand-coded audit of the reflections that explains why a richer signal did not help
(Section~\ref{sec:why}).

\section{Background and setup}

\textbf{How BFCL grades.} A task supplies the model with a set of simulated APIs
--- a file system, a trading bot, a travel booking service --- and a
conversation in which the user makes one request per turn. After each turn, two
gates must both pass. The \emph{state gate} replays the model's calls and
requires every public attribute of every involved simulator to equal the ground
truth's. The \emph{response gate} requires the ground truth's tool outputs to be
a subset of the model's.

This has a consequence that shapes several results below. Extra
\emph{read-only} calls are free: an agent may call \texttt{ls} or
\texttt{get\_balance} as often as it likes, since neither gate penalises
additional output or unchanged state. Extra \emph{state-changing} calls are not
--- an agent that satisfies every request but also alters something the ground
truth left alone fails the state gate (Section~\ref{sec:why} gives a case).

\textbf{The four categories.} Besides \texttt{miss\_param} and \texttt{miss\_func}
above, base is ordinary multi-turn tool use and \texttt{long\_context} pads the environment
so the relevant details are buried. The two abstention categories pull in opposite directions:
\texttt{miss\_param} penalises an agent that acts when it should have asked, while \texttt{miss\_func}
penalises one that hesitates when it should have acted. An agent that shifts its general willingness
to act will therefore gain in one category and lose in the other, a pattern that recurs throughout the results.

\textbf{Setup.} One model throughout, \texttt{qwen3-next-80b-a3b-instruct},
served over a university endpoint. A fixed subset of 200 tasks
(50 per category, stratified by which APIs a task involves, frozen to a
committed file and shared by every arm). \texttt{temperature=0}, one run per
task, an 8192-token cap that never binds, and a 30-step-per-turn guard that
fires on 1--4 tasks per arm. All code, prompts, trajectory logs and per-task
results are public \cite{kamilov2026code}.

Three fairness rules hold everywhere, because breaking any one of them would
make the comparison meaningless. Every arm emits \emph{actions} through the same
native function-calling channel, so no arm is penalised by a text-parsing
failure the others do not face; only the added \emph{thinking} is free text.
Every arm is zero-shot. Every arm shares one identical task-instruction core in
its system prompt, and differs only by its reasoning or reflection scaffold
(Appendices~\ref{app:sysprompt} and~\ref{app:react}).

\section{Method}

Six conditions run on the same 200 tasks. \textbf{Baseline} (plain function calling)
and \textbf{ReAct} do not retry. The other four share a single retry loop and differ
only in what fills its reflection slot: a fixed sham paragraph (\textbf{Sham}),
a shorter one with the advice removed (\textbf{Sham-lite}) (Appendix ~\ref{app:shams}),
the model's own lesson about the failure (\textbf{Reflexion}), or that lesson plus a record
of what the agent's own actions changed (\textbf{Reflexion+State}). Table~\ref{tab:arms} states each precisely.

\subsection{The shared retry loop}

Reflexion retries up to $k=3$ times. Retries are \textbf{fresh-episode}: the
failed conversation is discarded, the simulated environment is reset to its
initial configuration, and the retry is a new conversation that opens
with a short user message before turn~0 saying that a previous, separate attempt
failed, carrying the failure signal and the lesson (Appendix~\ref{app:preamble}).
Discarding the transcript is what makes the sham control meaningful: were the
failed conversation left in context, an agent given an empty lesson could
re-read it and diagnose the failure itself, silently acquiring the information
the sham is supposed to lack.

Attempt~1 of every retry arm reuses the recorded baseline run rather than executing
a fresh one. A task the baseline passed is carried through as a pass, with no reflection
written and no retry run; only the 115 tasks it failed enter the retry loop. Reflexion@1 is
therefore the baseline by construction, and all four retry arms retry an identical failure set.

The failure signal is an oracle: the grader tells the agent that it failed. Baseline and ReAct
never receive it, so any comparison between them and a retry arm carries that advantage;
all four retry arms receive the identical signal, so it cancels in the contrasts this report relies on.
It is also \textbf{sanitised}, because the grader's raw output contains the expected answer
--- the agent is told only which API class mismatched, or on which turn, and never any value (Appendix~\ref{app:signals}).

\subsection{Why a sham control is necessary}

Because the endpoint is not reproducible, a second attempt can succeed even
when the lesson contributes nothing: 7.8\% of failed tasks pass on an identical
re-run. The gap from @1 to @$k$ therefore confounds two quantities:

\[
\text{@1}\rightarrow\text{@}k \;=\; \underbrace{\text{value of the lesson}}_{\text{the quantity of interest}} \;+\; \underbrace{\text{effect of re-attempting}}_{\text{confound}}
\]

One measurement cannot separate two unknowns. The remedy is a control arm identical in every way except
that the lesson slot holds a fixed constant --- a \emph{sham} reflection, with
the same three headings and comparable length, written by no one
(Appendix~\ref{app:shams}). Then
\textbf{Reflexion@$k$ $-$ Sham@$k$} isolates the reflection content, and
\textbf{Sham@$k$ $-$ @1} absorbs the luck and the framing.
\textbf{Sham-lite} bounds how much of that comes from the advice rather than from retrying.

\subsection{The state signal}

The contribution arm, \textbf{Reflexion+State}, is identical to Reflexion except
that the reflection call also sees the agent's \textbf{own per-turn state
timeline}: what each turn changed, relative to the turn before
(Appendices~\ref{app:reflection} and~\ref{app:timeline}).\footnote{Arms are
named \texttt{baseline}, \texttt{react}, \texttt{reflexion},
\texttt{blind\_retry}, \texttt{blind\_retry\_lite} and
\texttt{richer\_reflexion\_turnwise} in the code and result files.}

This is not oracle-derived, which is what keeps it within the leak constraint.
The agent's own recorded calls are replayed from the task's initial
configuration, so nothing about the expected answer enters. This added timeline block is produced
once per reflection call, 218 times over the run: a genuine per-turn timeline in 175 cases; in 28,
a statement that the agent changed nothing --- itself diagnostic, since failing to act
is the characteristic error in the abstention categories; and in 15, all long\_context,
a placeholder, because the timeline exceeded a 3500-character cap set to stop a large filesystem swamping the prompt.

\subsection{Analysis}

All arms share the frozen subset, so every comparison is \textbf{paired per
task}. I report McNemar's exact test~\cite{mcnemar1947}, which discards the
tasks on which two arms agree and asks whether the remaining disagreements are
too one-sided to be attributed to chance, plus 20{,}000-sample bootstrap
confidence intervals over tasks.

The \textbf{minimum detectable effect} follows from the noise: with roughly 22
disagreements per comparison, significance needs about a 16--6 split, which is a
true effect of $\ge 5$~pp pooled and 8--10~pp within a single 50-task category.
I state this up front because most results below are smaller than that. Per
category, results are descriptive only.

\section{Results}

\subsection{The noise floor, and how much of it is the benchmark}
\label{sec:noise}

Two identical baseline runs disagree on \textbf{15 of 200 tasks}. Of the 115
tasks that failed the first run, 9 (7.8\%) pass on the second by chance alone.
This is the ruler for everything else.

I hand-labelled all 15 flips against the ground truth. Ten are genuine model
errors. \textbf{The remaining five are not instability at all: the agent's
behaviour was defensible and the grader rejected it.} Two examples show their
character. In \texttt{base\_37} the agent wrote a file with \texttt{echo},
which creates the file when it is absent; the ground truth first called
\texttt{touch} and then \texttt{echo}, both achieve the requested effect, but only
one satisfies the response gate. The final state is identical and the task fails anyway.
In \texttt{miss\_param\_119} the agent
verified that a user was logged in before closing their support ticket --- sound
practice in any real system --- but \texttt{close\_ticket} requires no login in
the simulator, and the extra check left state the ground truth never produced.
Table~\ref{tab:confounders} lists all five with their mechanisms.

\begin{table}[t]
\centering\footnotesize
\caption{The five benchmark artefacts among the 15 run-to-run flips. These are
grading brittleness, not model stochasticity.}
\label{tab:confounders}
\begin{tabular}{@{}ll@{}}
\toprule
\textbf{Task} & \textbf{Mechanism} \\
\midrule
\texttt{base\_117}       & Solved via a synonymous tool the ground truth did not choose \\
\texttt{base\_140}       & Ground truth demands exact free text the request underdetermines \\
\texttt{base\_37}        & Correct final state; skipped a redundant GT call returning \texttt{None} \\
\texttt{miss\_param\_119}& Real-world-sound auth check penalised; the tool needs no login \\
\texttt{miss\_param\_145}& Two symmetric login tools return contradictory status \\
\bottomrule
\end{tabular}
\end{table}

One category is weaker than it looks. In \texttt{miss\_func}, roughly half the
tasks never call the revealed tool at all (24/50 for the baseline, 21/50 for
ReAct). Reading the transcripts, the model denies the tool exists --- ``I don't
see any new tools added'' --- while the request verifiably contained it. That is
a failure of \emph{perception} of the tool list, which happens before any
reasoning loop could help. Expect reflection to be flat here, and it largely is.

\subsection{ReAct: a manipulation with nowhere to go}

ReAct scores 86/200 against the baseline's 85, a difference of $+0.5$~pp
(12 wins, 11 losses, $p=1.00$) --- indistinguishable from the two baseline runs'
disagreement with each other. The manipulation check explains why. With no
scaffold at all, the baseline already writes free-text reasoning before
\textbf{80.2\%} of its tool calls. ReAct moves that to \textbf{83.0\%}
($+3.1$~pp, CI $[-0.3,+6.4]$, sign test $p=0.033$) and adds about 20 characters
per action ($p=0.003$). The literal \texttt{Thought:} label appears on only
21.6\% of action steps. So the scaffold did shift behaviour slightly, but the
behaviour it targets was already near ceiling, and the pass rate did not move.

To be precise, what was tested is zero-shot ReAct-style thought injection over
native function calling --- not ReAct as published, which was demonstrated mainly
on a pre-trained (not instruction-tuned) model, with few-shot exemplars and actions
parsed out of text --- two differences forced by the fairness rules above,
which require every arm to be zero-shot and to act through the native function-calling channel.

\subsection{The ladder: only the first step matters}
\label{subsec:ladder}

Table~\ref{tab:arms} summarises the Reflexion experiments: the four retry arms,
together with the two non-retrying conditions they are measured against. Signal
content increases at every rung and the score does not. The four retry arms span
three tasks (100 to 103) --- the same spread as two identical runs of the
Baseline, which scored 85 and 88.

\begin{table}[t]
\centering\footnotesize
\caption{The six conditions, ordered by how much the retry is told. The lower
four are the Reflexion experiments; Baseline and ReAct do not retry and are
shown for reference. pass@3 is out of 200 tasks.}
\label{tab:arms}
\begin{tabular}{@{}llc@{}}
\toprule
\textbf{Arm} & \textbf{What it adds} & \textbf{pass@3} \\
\midrule
Baseline           & native function calling, no scaffold      & 85 \\
ReAct              & a thought before each action              & 86 \\
\midrule
Sham-lite          & retry + sanitised signal + a bland lesson & 102 \\
Sham               & \quad + generic corrective advice         & 103 \\
Reflexion          & \quad + a lesson written from the failure & 100 \\
Reflexion+State    & \quad + the agent's own state timeline    & 100 \\
\bottomrule
\end{tabular}
\end{table}
The one large effect is the retry itself: Sham reaches 103 against the 85
every arm shares at @1 (first-attempt --- seeded from Baseline), a gain of
$\mathbf{+9.0}$~pp (18 wins, 0 losses, $p<10^{-4}$) --- necessarily one-sided,
since a retry fires only on a failure and inherits every pass. After that,
nothing moves.

\begin{itemize}\itemsep2pt
\item \textbf{Real reflection loses to the sham.} Reflexion $-$ Sham is
  $\mathbf{-1.5}$~pp, CI $[-5.5,+2.5]$, 6 wins to 9 losses, $p=0.607$.
\item \textbf{The sham's advice is inert.} Sham $-$ Sham-lite is $+0.5$~pp
    (2 to 1, $p=1.00$), with three of four categories identical task-for-task. The
    sham arms flip 12.2\% of failures on the first retry against 7.8\% for a plain
    re-run, so that 4.4~pp belongs to the retry framing and the sanitised signal,
    not to advice.
\item \textbf{The richer signal changes nothing.} Reflexion+State $-$ Reflexion
  is $\mathbf{+0.0}$~pp, CI $[-3.5,+3.5]$, 7 to 7, $p=1.00$.
\end{itemize}

This is the result the control bought. Without a sham arm, the honest-looking
headline would have been ``Reflexion $+7.5$~pp, $p<10^{-4}$'', and every point
of it would have been retry luck.

\subsection{Which failures, not how many}
\label{sec:which}

The pooled zero is not a flat null. It is two opposite effects cancelling
(Table~\ref{tab:categories}).

\begin{table}[t]
\centering\footnotesize
\caption{Pass counts per category (of 50). The pooled null hides
opposite-signed category effects. Descriptive only: per category the detectable
effect is 8--10~pp.}
\label{tab:categories}
\begin{tabular}{@{}lcccrr@{}}
\toprule
& \multicolumn{3}{c}{\textbf{pass@3}} & \multicolumn{2}{c}{\textbf{difference}} \\
\cmidrule(lr){2-4}\cmidrule(lr){5-6}
\textbf{Category} & Sham & Reflexion & +State & Refl.$-$Sham & +State$-$Refl. \\
\midrule
\texttt{base}          & 38 & 37 & 37 & $-2$\,pp  & $\pm0$\,pp \\
\texttt{miss\_func}    & 10 & \textbf{14} & 10 & $\mathbf{+8}$\,pp & $\mathbf{-8}$\,pp \\
\texttt{miss\_param}   & \textbf{20} & 15 & \textbf{19} & $\mathbf{-10}$\,pp & $\mathbf{+8}$\,pp \\
\texttt{long\_context} & 35 & 34 & 34 & $-2$\,pp  & $\pm0$\,pp \\
\bottomrule
\end{tabular}
\end{table}

Reflection reliably moves one thing: the agent's threshold for \emph{acting}
versus \emph{asking}, since its lessons converge on generic rules about when to
act and when to check with the user. In \texttt{miss\_func}, where the dominant
failure is under-acting, that push helps ($+8$~pp). In \texttt{miss\_param},
which punishes hesitancy, the same push over-corrects into caution and hurts
($-10$~pp); lessons such as ``always confirm the user ID before acting'' are
exactly wrong there.

Adding the state signal moves the threshold back, in both directions at once:
\texttt{miss\_param} recovers ($15\to19$, back to the sham's level) at the exact
cost of \texttt{miss\_func} ($14\to10$, also back to the sham's level). The
behaviour confirms it. Tool calls per retry in \texttt{miss\_param} run 8.00 for
Reflexion, 8.87 for Reflexion+State and 9.33 for the Sham --- the arms are
ordered by how willing they are to act. So whether reflection helps is a property
of the task category, not of the reflection. Averaged over four categories it cancels.

\subsection{The ceiling: a small fixed pool}
\label{sec:ceiling}

There is a simpler reason all four retry arms land within three tasks of each
other: re-attempting does not act uniformly on the 115 failures. A small pool of
borderline tasks flips readily and the large majority are locked. Across five
independent re-roll conditions only \textbf{28 of the 115 failures are ever
recovered by anything}, and each arm recovers 15--18 of those same 28 (Sham 18,
Sham-lite 17, Reflexion 15, Reflexion+State 15). The other 87 never flip for
anyone; they are the \texttt{miss\_func} tool denials and the grading artefacts
of Section~\ref{sec:noise}. Reflection therefore reallocates within a fixed pool
rather than enlarging it: Reflexion and the Sham overlap on only 38\% of the
tasks they convert, and the split is category-structured.

\subsection{Cost}

The retry arms cost about $2.1\times$ the Baseline's input tokens (29.5M against 13.8M),
and that multiple is almost entirely the extra attempts. Within Reflexion the reflection
mechanism accounts for roughly 5\% of the tokens --- 2\% to generate the lessons, and about
3\% more to carry them, since a lesson sits in the retry's opening message and is re-sent with
every call in that attempt. The state timeline adds a further 0.1\%. So the expensive
component of Reflexion is the retry, which worked, and the cheap component is the reflection, which did not.

\section{Why the richer signal did not help}
\label{sec:why}

\subsection{It changed salience, not information}

The state timeline is largely something the agent already had. Across all 175 real
timelines (654 attribute lines), \textbf{47.6\% of the values appear verbatim in
the transcript the reflection step is already reading}, and most of the rest is
derivable from tool results in that same transcript. The timeline block restructures and
aggregates; it does not inform. This is a lexical test and should be read as
one, but it sets the scope of the null: what remains untested is a signal
genuinely \emph{new} to the agent, such as an external validator or an
environment probe it could not otherwise observe.

\subsection{A hand-coded audit of the reflections}

The two reflection arms disagree on 14 tasks. Their attempt-1 reflections form
an unusually clean comparison, because the shared seeding makes the transcript,
the signal, the prompt template and the temperature identical, leaving the state
block as the only difference. I read all 28 against the ground truth and coded
them on two axes: whether the \emph{diagnosis} names the real cause, and
whether following the \emph{lesson} would help (Table~\ref{tab:audit}).
The two arms are level, with Reflexion if anything marginally ahead on diagnosis.
The block changed the conclusion on 6 of the 14 tasks, but symmetrically ---
improving the diagnosis twice and degrading it twice --- and it went unused
about as often as it helped: in \texttt{base\_94} the agent released a parking
brake the ground truth left engaged, and the block lists that exact change while
the lesson blames tyre pressure instead.

\begin{table}[t]
\centering\footnotesize
\caption{Hand-coded quality of the 14 discordant attempt-1 reflection pairs. No
bucket differs by more than two tasks.}
\label{tab:audit}
\begin{tabular}{@{}lccc c ccc@{}}
\toprule
& \multicolumn{3}{c}{\textbf{Diagnosis}} & & \multicolumn{3}{c}{\textbf{Lesson}} \\
\cmidrule(lr){2-4}\cmidrule(lr){6-8}
\textbf{Arm} & correct & partial & wrong & & correct & partial & wrong \\
\midrule
Reflexion        & \textbf{6} & 3 & 5 & & \textbf{6} & 2 & 6 \\
Reflexion+State  & 5 & 4 & 5 & & 5 & 4 & 5 \\
\bottomrule
\end{tabular}
\end{table}

\subsection{Two tasks that show the whole mechanism}

\textbf{Where it worked.} In \texttt{miss\_param\_32} the ground truth requires
writing \texttt{2.0} into a file. Both arms diagnose the failure correctly: the
agent wrote an explanatory sentence into the file instead of the bare number.
They differ on the fix. Reflexion \emph{guessed} the format --- ``output only
the precise numeric value (e.g.\ \texttt{2.0000})'' --- wrote
\texttt{2.0000} on its retry, failed again, and doubled down. It lost because of
its own lesson. Reflexion+State was shown the file it had actually written and
prescribed ``write only the raw number (e.g.\ \texttt{2.0})''. It passed. Note
what the block contributed: not understanding, which both arms had, but a
concrete anchor that stopped the fix drifting into invention.

\textbf{Where it backfired.} In \texttt{miss\_param\_169} Reflexion diagnosed the
failure exactly right --- the user had not yet given the travel date or class, so
the agent should have asked. Reflexion+State did not. Its state block showed that its own booking attempt
had changed a credit-card record, and it concluded that ``the system required
\texttt{client\_id}, \texttt{client\_secret}, and other authentication details''.
Those parameters are real. They belong to \texttt{authenticate\_travel}, one of the 27 tools the agent held.
But booking does not need them --- only the \texttt{access\_token} the user had already supplied.
So the block did not make the model invent a requirement. It prioritised the wrong tool,
and the lesson sent the retry hunting for credentials instead of asking for the missing details.
Same mechanism, opposite directions --- that is what a $+0.0$~pp result is made of.

\section{Limitations}
\enlargethispage{3\baselineskip}

\textbf{One model.} All results are specific to
\texttt{qwen3-next-80b-a3b-instruct}; the ReAct saturation finding in particular
is a property of an instruction-tuned model that already narrates.

\textbf{Two of three loop positions.} I test before-acting and after-episode.
Per-step reflection was cut for time, and it is the most promising untested
position precisely because of this result: episodic reflection shifts the
act/ask threshold \emph{globally}, which is why opposite category effects
cancel, whereas a per-step variant could apply the correction locally.

\textbf{No pass@$k$ figure is deployment-realistic.} A retry fires only when
the grader reports a failure, and a deployed agent has no such signal.
It also restarts from a clean environment: an agent that has already booked
a flight or deleted a file cannot take that back. The setup is realistic as a
data-generation procedure rather than as an architecture. A sandbox supplies exactly
what deployment lacks --- a failure oracle and an undo --- so an agent
can attempt a task repeatedly and the trajectories that pass can be used to train the model.
Yao et al. apply the same generate-filter-train pattern, fine-tuning on
3{,}000 ReAct trajectories selected for task success~\cite{yao2023react}.

\textbf{The richer signal manipulated salience, not information}
(Section~\ref{sec:why}). The state timeline is close to everything BFCL permits
without leaking the answer: it exposes no test suite, no validator, leaving only
the transcript, a sanitised verdict and the agent's own effect on the environment.
A setting offering more --- the stack traces of code generation, or an environment
the agent may probe --- would give it information it could not otherwise obtain.
That condition could not be constructed here.

\textbf{Benchmark validity.} Five of the fifteen run-to-run flips are grading
artefacts, and \texttt{miss\_func} substantially measures tool-list perception
rather than reasoning (Section~\ref{sec:noise}).

\section{Conclusion}

On this benchmark and this model, \emph{where} reflection sits in the loop
proved less informative than \emph{what it is told}: even a substantially richer
input did not change how many failures were recovered, only which ones. The
number is capped by a small pool of borderline tasks any retry can reach, and
reflection's observable effect is to shift the agent's act/ask threshold,
helping in one category as much as it hurts in another.

The methodological lesson generalises: when the system under test is not
reproducible and re-attempting is part of the architecture, an uninformative
control arm is necessary rather than optional. Measured against its own first attempt,
Reflexion gains $7.5$~pp at $p<10^{-4}$. Measured against a sham that says
nothing, it gains nothing. Only the second comparison controls for re-attempting.

\bibliographystyle{abbrv}
\bibliography{references}

\clearpage
\appendix
\section*{Appendix: prompts and signals}
\addcontentsline{toc}{section}{Appendix}
\renewcommand{\thesubsection}{A.\arabic{subsection}}
\setcounter{subsection}{0}

\noindent
Every prompt below is reproduced verbatim from the committed source
\cite{kamilov2026code}. Braces such
as \texttt{\{signal\}} are substitution placeholders filled at run time. Together
these make the report's central claim checkable: at each rung of the ladder the
arms differ in exactly one place, and that place is printed here.

\subsection{Shared system prompt core (Baseline)}
\label{app:sysprompt}
Used unchanged by \emph{every} arm. The Baseline's system prompt is exactly this
text and nothing else; other arms append a scaffold to it.
\begin{lstlisting}
You are an agent that completes user tasks by calling the provided tools. Call tools whenever they are needed to satisfy the user's request.
If none of the available tools can fulfil the request, say so instead of guessing. If the request is missing a value that a tool requires, ask the user for it.
Keep calling tools until the current request is fully handled. When you have nothing left to call, reply in natural language summarising what you did — a message with no tool call ends the current turn.
\end{lstlisting}

\subsection{ReAct scaffold}
\label{app:react}
Appended to the core above, after a blank line. This is the entire ReAct
manipulation.
\begin{lstlisting}
Work in an explicit reason-act loop. Every time you call a tool, begin that SAME message with the literal label 'Thought:' followed by one or two sentences — what you know, what you still need, and which tool to call next — then make the tool call(s) in that same message.
A Thought with no tool call attached ends the current turn. After each tool result, any further tool call again begins with a 'Thought:' label in that same message.
\end{lstlisting}

\subsection{The reflection prompt, and the state variant}
\label{app:reflection}
Reflexion's reflection call. It is a single LLM call with no tools, at the same
temperature and token budget as the action loop.
\begin{lstlisting}
Below is the full transcript of your previous attempt at a multi-turn tool-calling task. The attempt was graded as FAILED.

{attempt_text}

Grader failure signal: {signal}

Reflect on this failed attempt. Reply in exactly this format:
What I tried: <one or two sentences>
What went wrong: <one or two sentences — your best diagnosis, using the failure signal>
Lesson: <one or two sentences of concrete, actionable guidance>

Whoever reads this reflection later will not see this transcript — only the reflection itself. Make the Lesson self-contained.
\end{lstlisting}

\noindent
Reflexion+State uses a prompt that is \textbf{byte-identical to the above}
except for three inserted lines after the failure signal (551 characters against
649). This is the whole of the contribution:
\begin{lstlisting}
+
+For reference, here is how your own actions changed the environment state, turn by turn:
+{state}
\end{lstlisting}

\subsection{The two sham reflections}
\label{app:shams}
The Sham arm's lesson slot holds this fixed constant. It is written by no one and
contains nothing about the failure, but it is form-matched to
Appendix~\ref{app:reflection}'s output: the same three headings, comparable
length.
\begin{lstlisting}
What I tried: I attempted the task using the tools provided.
What went wrong: Some part of my approach did not match what the task required.
Lesson: I should re-read the request carefully, verify preconditions before acting, and double-check tool arguments.
\end{lstlisting}

\noindent
Sham-lite is identical except that the \texttt{Lesson} line drops its two
imperative clauses, giving \texttt{"Lesson: I should re-read the request
carefully."} That is the entire difference between the two control arms: 257
characters against 188, and they score 103 and 102 out of 200.

\subsection{The retry preamble}
\label{app:preamble}
The user message every retry opens with, before turn~0. The template is
\textbf{identical in all four retry arms}; only the reflection text differs. One
\texttt{signal}/\texttt{reflection} pair is appended per failed attempt.
Example, attempt~2 of a Sham run:
\begin{lstlisting}
Note: a previous, separate attempt at this task failed.
Attempt 1 failure signal: After your attempt, the TicketAPI environment was left in a state that does not match what the task required.
Attempt 1 reflection: What I tried: I attempted the task using the tools provided.
What went wrong: Some part of my approach did not match what the task required.
Lesson: I should re-read the request carefully, verify preconditions before acting, and double-check tool arguments.
The task now restarts from the beginning.
\end{lstlisting}

\subsection{Sanitized failure signals}
\label{app:signals}
The grader's raw verdict leaks the expected answer, so it is mapped to one of
four templates carrying only an environment class or a turn index. These are the
complete set; no value from the ground truth ever reaches the agent.
\begin{lstlisting}
[state gate]     After your attempt, the {ClassName} environment was left in a state that does not match what the task required.

[response gate]  By the end of user turn {n} (counting from 1), the output of one or more required tool calls was missing: a call the task required was never made, or was made with different arguments.

[empty turn]     You ended user turn {n} (counting from 1) without making any tool calls, but that turn required action.

[irrelevance]    You made tool calls in user turn {n} (counting from 1), but the correct behavior there was to make none (e.g. ask the user for missing information, or say that no available tool fits).
\end{lstlisting}

\noindent
The state gate carries no turn index: the checker does not record which turn
produced the mismatch, so those signals are class-level only.

\subsection{An example state timeline}
\label{app:timeline}
The block Reflexion+State inserts at \texttt{\{state\}}, here the real one from
\texttt{miss\_param\_32} (Section~\ref{sec:why}): it shows the agent the file it
actually wrote, which let it correct \texttt{2.0000} to \texttt{2.0}. Turns
changing nothing are listed as such --- the under-acting signal.
\begin{lstlisting}
Turn 1: no state changes
Turn 2: no state changes
Turn 3:
  GorillaFileSystem:
    root = "<Directory: project, Parent: None, Contents: {'Spring2023Draft': <<File: Spring2023Draft, Content: These are the notes for Spring 2023.>>, 'PastSeasons': <Directory: PastSeasons, Parent: project, Contents: {}>, 'result.txt': <<File: result.txt, Content: The logarithm of 36 to base 6 is: 2.0>>}>"
\end{lstlisting}

\end{document}
